\section{Conclusion}\label{sec:conclusion}
\noindent We presented a five-step framework for evaluating IoMT intrusion detection at the network level, from arithmetic audit to simulation, and applied it to three published systems and four public datasets. The audit confirmed metric-definition errors in one paper and a split drawn after oversampling in two, and showed experimentally that the ordering alone can add 0.29 to F1. Duplicate rates depend on the feature space retained, from under 1\,\% to 99\,\% on the same file. In an ns-3 fog network, rules with identical detection rates quarantined between 45\,\% and 100\,\% of legitimate wearables in the baseline, and the published rule all of them in every scenario; a congestion-coupled rule latched off under load, topology-driven concentration and a measured tail inference cost, and three mitigations each restored it in every run; and threshold adaptation proved bounded by the information it acts on, with a low-false-positive classifier gate moving that bound. On a Raspberry Pi 4 the trained ensemble reproduces its evaluated predictions exactly and costs 3.0\,ms per window, enough to switch a congestion-coupled rule off at the Review~2 attack scale. Three items remain and are scripted: the reader study, the 20-partition IoTID20 tables and the public release.

\section*{Reproducibility}\label{sec:repro}
\noindent The ns-3 simulator, the Review~2 program used for the regression check, every experiment definition, all raw runs (one line per run, written before aggregation), the dataset-leakage and multi-partition scripts, the audit extraction with source locations, the hardware harness, the reader-study materials and the manuscript build are released under the MIT licence, with a citation file and archival metadata for a Zenodo DOI. Datasets are not redistributed; their sources, checksums and row counts are listed with the code.

\section*{Acknowledgment}
\noindent The authors thank the course faculty for the Review~3 roadmap, the maintainers of the datasets used here, and the authors of \cite{berguiga2025}, \cite{hosain2025} and \cite{tarish2026}, whose explicit reporting made the audit possible.
